% ADDED figure (2026-11-17): annihilation and scattering side by side with the crossed legs marked (after PHY 513 Lecture 21 slide 6 and PS p. 153); generated by qft-C/tools/l21/figs/make_figs.sh
\tikzset{
  fdl/.style={line width=0.8pt, black},
  fdx/.style={circle, fill=black, inner sep=0pt, minimum size=3.6pt},
  fdv/.style={circle, fill=figB, draw=figB, inner sep=0pt, minimum size=5pt},
  fdlab/.style={font=\scriptsize, inner sep=2.6pt},
  fdmom/.style={->, >=Stealth, figO, line width=0.6pt},
}
\tikzset{
  fda/.style={fdl, postaction={decorate}, decoration={markings, mark=at position 0.55 with {\arrow{Stealth[length=5pt]}}}},
  fdS/.style={font=\small},
}
\tikzset{
  fline/.style={line width=0.8pt, black, postaction={decorate}, decoration={markings, mark=at position 0.56 with {\arrow{Stealth[length=5.5pt, width=4.5pt]}}}},
  sline/.style={line width=0.8pt, black, dash pattern=on 3.5pt off 2.2pt},
  blob/.style={draw, line width=0.8pt, fill=figB!12},
}
\tikzset{
  pline/.style={line width=0.8pt, black, decorate, decoration={snake, amplitude=1.6pt, segment length=5.2pt}},
}
\begin{tikzpicture}[font=\small]
  % (a) annihilation, time to the right
  \begin{scope}
    \coordinate (v1) at (0,0); \coordinate (v2) at (2.2,0);
    \draw[fline] (-1.4,-1.4)--(v1);
    \draw[fline, figR] (v1)--(-1.4,1.4);
    \draw[fline] (v2)--(3.6,-1.4);
    \draw[fline, figR] (3.6,1.4)--(v2);
    \draw[pline] (v1)--(v2); \node[fdv] at (v1) {}; \node[fdv] at (v2) {};
    \node[fdlab, left] at (-1.4,-1.4) {$e^-\ (p)$};
    \node[fdlab, left, figR] at (-1.4,1.4) {$e^+\ (p')$};
    \node[fdlab, right] at (3.6,-1.4) {$\mu^-\ (k)$};
    \node[fdlab, right, figR] at (3.6,1.4) {$\mu^+\ (k')$};
    \node[fdlab, above] at (1.1,0.12) {$q = p + p'$};
    \draw[->, gray] (0.2,-1.95)--(2.0,-1.95); \node[font=\scriptsize, gray, right] at (2.0,-1.95) {time};
    \node[fdlab] at (1.1,2.25) {(a) $e^-e^+ \to \mu^-\mu^+$};
    \node[fdlab] at (1.1,-2.75) {$q^2 = E_{\rm cm}^2 > 0$};
  \end{scope}
  % substitutions
  \node[align=left, font=\footnotesize, anchor=west] at (5.15,0.15) {$p \to p_1$\\[3pt] \textcolor{figR}{$p' \to -p_1'$}\\[3pt] $k \to p_2'$\\[3pt] \textcolor{figR}{$k' \to -p_2$}};
  \draw[->, line width=0.7pt] (5.2,-1.35)--(6.7,-1.35);
  % (b) scattering, time upward
  \begin{scope}[shift={(9.3,0)}]
    \coordinate (v1) at (0,0); \coordinate (v2) at (2.2,0);
    \draw[fline] (-1.4,-1.4)--(v1);
    \draw[fline, figR] (v1)--(-1.4,1.4);
    \draw[fline, figR] (3.6,-1.4)--(v2);
    \draw[fline] (v2)--(3.6,1.4);
    \draw[pline] (v1)--(v2); \node[fdv] at (v1) {}; \node[fdv] at (v2) {};
    \node[fdlab, left] at (-1.4,-1.4) {$e^-\ (p_1)$};
    \node[fdlab, left, figR] at (-1.4,1.4) {$e^-\ (p_1')$};
    \node[fdlab, right, figR] at (3.6,-1.4) {$\mu^-\ (p_2)$};
    \node[fdlab, right] at (3.6,1.4) {$\mu^-\ (p_2')$};
    \node[fdlab, above] at (1.1,0.12) {$q = p_1 - p_1'$};
    \draw[->, gray] (5.6,-0.9)--(5.6,0.9); \node[font=\scriptsize, gray, right] at (5.6,0) {time};
    \node[fdlab] at (1.1,2.25) {(b) $e^-\mu^- \to e^-\mu^-$};
    \node[fdlab] at (1.1,-2.75) {$q^2 \le 0$};
  \end{scope}
  \node[font=\footnotesize, align=center] at (6.6,-3.6) {black: the same particle, renamed \qquad \textcolor{figR}{red: crossed (incoming $e^+$ $\leftrightarrow$ outgoing $e^-$; outgoing $\mu^+$ $\leftrightarrow$ incoming $\mu^-$)}};
\end{tikzpicture}
